\DSource{0131}{42}
\hypertarget{AB01-PDF0131-P01}{}
\DLine{AB01-PDF0131-body-L001}{}{}{et observationem Hipparchi idem fere tempus intercessisse quod inter Hipparchi et Ptolemaei;}
\DLine{AB01-PDF0131-body-L002}{}{}{fuit enim observatio 286 (1) annis ante Hipparchum.}
\hypertarget{AB01-PDF0131-P02}{}
\DLine{AB01-PDF0131-body-L003}{}{}{\Indent Nos ipsi in urbe ar-Raqqah observavimus; et una nostrarum observationum autumna-}
\DLine{AB01-PDF0131-body-L004}{}{}{rum, qua confidimus et de cuius veritate, ut ex instrumentis apparuit, haud dubitamus, est}
\DLine{AB01-PDF0131-body-L005}{5}{}{observatio quae 743 annis fuit post observationem autumnam Ptolemaei iam memoratam. Ea}
\DLine{AB01-PDF0131-body-L006}{}{}{invenimus Solem per punctum aequinoctii autumni transivisse, anno 1194 aerae Dhū ’l-qar-}
\DLine{AB01-PDF0131-body-L007}{}{}{nayn, idest 1206 post Alexandrum mortuum, ante Solis ortum diei XIX mensis Romani Aylūl (2),}
\DLine{AB01-PDF0131-body-L008}{}{}{scilicet VIII diei mensis Coptici Pachōn, quatuor circiter horis et $\qfrac{1}{2}+\qfrac{1}{4}$ ({\itshape c}). Et quia circulus}
\DLine{AB01-PDF0131-body-L009}{}{}{meridianus Alexandriae circulum meridianum urbis ar-Raqqah $\qfrac{2}{3}$ horae aequinoctialis circiter}
\DLine{AB01-PDF0131-body-L010}{10}{}{praecedit (3), inter observationem nostram et observationem autumnam Ptolemaei 743 anni}
\DLine{AB01-PDF0131-body-L011}{}{}{Aegyptii, 178 dies, et $\qfrac{1}{2}+\qfrac{1}{4}$ diei, detractis $\qfrac{2}{5}$ horae circiter (4), intercesserunt, loco 185}
\DLine{AB01-PDF0131-body-L012}{}{}{dierum $\qfrac{1}{2}+\qfrac{1}{4}$ qui debuissent in his annis colligi si quartae partes diei integrae fuissent. Si}
\DLine{AB01-PDF0131-body-L013}{}{}{hos 7 dies et $\qfrac{2}{5}$ horae, quibus tempus observationis tempus quartae partis diei 365 dies supe-}
\DLine{AB01-PDF0131-body-L014}{}{}{rantis praecessit, per 743 annos qui inter duas observationes illas fuerunt dividimus, invenimus}
\DLine{AB01-PDF0131-body-L015}{15}{}{portionem unius anni esse $3^\circ24'$ ex iis $360^\circ$ qui revolutionem diei et noctis perficiunt. Quam}
\DLine{AB01-PDF0131-body-L016}{}{}{portionem si de tempore quartae partis diei, scilicet de $90^\circ$ detrahimus, quantitas augmenti}
\DLine{AB01-PDF0131-body-L017}{}{p. 64.}{super 365 dies integros remanet $86^\circ36'$; est igitur verum anni tempus $365^d14'26''$ (5) fere.}
\hypertarget{AB01-PDF0131-P03}{}
\DLine{AB01-PDF0131-body-L018}{}{}{\Indent Si $360^\circ$ circumferentiae caelestis sphaerae per tempus anni nuper inventum dividimus,}
\DLine{AB01-PDF0131-body-L019}{}{}{motum medium Solis in nychthemero reperimus $0^\circ59'8''20'''46^{\mathrm{IV}}56^{\mathrm{V}}14^{\mathrm{VI}}$ esse, motum au-}
\DLine{AB01-PDF0131-body-L020}{20}{}{tem in 30 diebus, qui mensem Aegyptium efficiunt, $29^\circ34'10''23'''28^{\mathrm{IV}}6^{\mathrm{V}}47^{\mathrm{VI}}$, et in 365}
\DLine{AB01-PDF0131-body-L021}{}{}{diebus anni Aegyptii $359^\circ45'46''25'''32^{\mathrm{IV}}2^{\mathrm{V}}31^{\mathrm{VI}}$ (6) circiter ({\itshape d}). Hos motus multiplicavimus}
\DLine{AB01-PDF0131-body-L022}{}{}{et in tabulis ad annos collectos et singulos, et ad menses, dies, horas, iuxta aeram Arabum}
\DLine{AB01-PDF0131-body-L023}{}{}{aeramque Romanorum descripsimus (7), ut facile quocumque harum aerarum tempore locus}
\DLine{AB01-PDF0131-body-L024}{}{}{Solis, secundum eius motum medium qui {\itshape medium Solis} appellatur, inveniri possit. Tempus}
\DLine{AB01-PDF0131-body-L025}{25}{}{anni a nobis inventum patet $2^\circ\qfrac{1}{5}$ minus esse quam tempus a Ptolemaeo memorato (8); qua}
\DLine{AB01-PDF0131-body-L026}{}{}{re motus Solis, iuxta calculum nostrum, motum a Ptolemaeo inventum in die $0^\circ0'0''3'''$}
\DLine{AB01-PDF0131-body-L027}{}{}{$33^{\mathrm{IV}}43^{\mathrm{V}}43^{\mathrm{VI}}$ (9) superat, et in anno Aegyptio, Deo volente, $0^\circ0'21''40'''10^{\mathrm{IV}}53^{\mathrm{V}}56^{\mathrm{VI}}$ fere (10).}
\par\vspace{6pt}\hrule height.25pt\vspace{5pt}
\noindent\begin{minipage}[t]{.48\textwidth}\vspace{0pt}\fontsize{9}{11.5}\selectfont\setlength{\GutterOffset}{0pt}
\hypertarget{AB01-PDF0131-N01}{}
\DLine{AB01-PDF0131-notes_left-L001}{}{}{\Indent (1) Codex 280, Plato 186.}
\hypertarget{AB01-PDF0131-N02}{}
\DLine{AB01-PDF0131-notes_left-L002}{}{}{\Indent (2) I. e. ante Solis ortum diei 19 Sept. 882;}
\DLine{AB01-PDF0131-notes_left-L003}{}{}{nam al-Battānī prima die mensis Septembris annos}
\DLine{AB01-PDF0131-notes_left-L004}{}{}{huius aerae incipere vult, cfr. cap. XXXII. --- Hunc}
\DLine{AB01-PDF0131-notes_left-L005}{}{}{locum laudat \Name{Ibn Yūnus} apud \Name{Caussin}, {\itshape Le li-}}
\DLine{AB01-PDF0131-notes_left-L006}{}{}{{\itshape vre de la grande table Hakémite} (Notices et ex-}
\DLine{AB01-PDF0131-notes_left-L007}{}{}{traits des mss. de la Bibl. Nation., t. VII, Paris 1804,}
\DLine{AB01-PDF0131-notes_left-L008}{}{}{p. 149).}
\hypertarget{AB01-PDF0131-N03}{}
\DLine{AB01-PDF0131-notes_left-L009}{}{}{\Indent (3) In suis tabulis geographicis, e communi}
\DLine{AB01-PDF0131-notes_left-L010}{}{}{Arabica traditione depromptis, al-Battānī urbi ar-}
\DLine{AB01-PDF0131-notes_left-L011}{}{}{-Raqqah $73^\circ15'$ long. tribuit, Alexandriae, Ptole-}
\DLine{AB01-PDF0131-notes_left-L012}{}{}{maeum sequens, $60^\circ30'$. Hinc differentia longitudi-}
\DLine{AB01-PDF0131-notes_left-L013}{}{}{nalis $12^\circ45'$ sequeretur, non $10^\circ$.}
\hypertarget{AB01-PDF0131-N04}{}
\DLine{AB01-PDF0131-notes_left-L014}{}{}{\Indent (4) I. e. 743 anni, $178^d\,17^h\,36^m$.}
\hypertarget{AB01-PDF0131-N05}{}
\DLine{AB01-PDF0131-notes_left-L015}{}{}{\Indent (5) I. e. $365^d\,5^h\,46^m\,24^s$. --- \Name{Ptolemaeus}, {\itshape Al-}}
\DLine{AB01-PDF0131-notes_left-L016}{}{}{{\itshape mag.} III, 2 (Halma t. I, p. 165), seu potius Hippar-}
\end{minipage}
\hfill\begin{minipage}[t]{.48\textwidth}\vspace{0pt}\fontsize{9}{11.5}\selectfont\setlength{\GutterOffset}{0pt}
\DLine{AB01-PDF0131-notes_right-L001}{}{}{chus, invenerat $365^d14'48''=365^d\,5^h\,55^m\,12^s$; astro-}
\DLine{AB01-PDF0131-notes_right-L002}{}{}{nomi nostri temporis $365^d\,5^h\,48^m\,47^s{,}33$ ponunt.}
\hypertarget{AB01-PDF0131-N06}{}
\DLine{AB01-PDF0131-notes_right-L003}{}{}{\Indent (6) Male quartae in cod. et Plat. $31^{\mathrm{IV}}$. Iam cor-}
\DLine{AB01-PDF0131-notes_right-L004}{}{}{rexit Ed. \Name{Halley}, {\itshape Emendationes ac notae in ve-}}
\DLine{AB01-PDF0131-notes_right-L005}{}{}{{\itshape tustas Albatênii observationes} (Philosophical Tran-}
\DLine{AB01-PDF0131-notes_right-L006}{}{}{sactions of the Royal Society, 1693, vol. XVII, nr. 204,}
\DLine{AB01-PDF0131-notes_right-L007}{}{}{p. 916.}
\hypertarget{AB01-PDF0131-N07}{}
\DLine{AB01-PDF0131-notes_right-L008}{}{}{\Indent (7) Vide tabulas in fol. 164,v.--166,v., et 186,v.-}
\DLine{AB01-PDF0131-notes_right-L009}{}{}{-189,r.}
\hypertarget{AB01-PDF0131-N08}{}
\DLine{AB01-PDF0131-notes_right-L010}{}{}{\Indent (8) Nam $8^m\,48^s\times15=528^s\times15=7920''=$}
\DLine{AB01-PDF0131-notes_right-L011}{}{}{$2^\circ12'=2^\circ\qfrac{1}{5}$.}
\hypertarget{AB01-PDF0131-N09}{}
\DLine{AB01-PDF0131-notes_right-L012}{}{}{\Indent (9) Codex male $30^{\mathrm{IV}}$ et $42^{\mathrm{VI}}$; \Name{Ptolemaeus} III,}
\DLine{AB01-PDF0131-notes_right-L013}{}{}{2 (ed. Halma, t. I, p. 166) enim habet $0^\circ59'8''17'''$}
\DLine{AB01-PDF0131-notes_right-L014}{}{}{$13^{\mathrm{IV}}12^{\mathrm{V}}31^{\mathrm{VI}}$.}
\hypertarget{AB01-PDF0131-N10}{}
\DLine{AB01-PDF0131-notes_right-L015}{}{}{\Indent (10) Male quintae apud Platonem $55^{\mathrm{V}}$ et in co-}
\DLine{AB01-PDF0131-notes_right-L016}{}{}{dice $50^{\mathrm{V}}$; \Name{Ptolemaeus} $359^\circ45'24''45'''21^{\mathrm{IV}}8^{\mathrm{V}}35^{\mathrm{VI}}$.}
\end{minipage}
